% ============================================================
% primitives.tex   —— 工具箱：画一页用的最小原语
%
% 【这是什么】框架的"工具"。它定义一页怎么搭起来，以及所有模板都会
%   调用的最小部件。**不含任何视觉设定**（那些在 styles/ 里），
%   **不含任何版式结构**（那些在 templates/ 里）。
%
% 【谁改它】几乎不改。只有框架本身新增能力时才动 —— 这是它和
%   config.tex 的本质区别：config.tex 是配置（可再生），
%   本文件是工具（扔了就没有框架了）。
%
% 【加载位置】config.tex 的**第一步**，在任何东西之前。
%   下面这些命令只**引用** \DeckRowSkip、\SlideAccent 这类名字，
%   定义时并不求值 —— 所以风格还没加载也完全没问题，
%   真正用到它们是在排版那一刻，那时风格早就位了。
%
%   为什么不能把本文件并进 config.tex：那会让"永不变的工具"和
%   "每个 deck 都不同的配置"住进同一个文件 —— 一个文件两个变动原因。
% ============================================================

% ============================================================
% 一、定义组件的命令
% ------------------------------------------------------------
% 模板里定义组件一律用 \DeckDefine，不要用 \newcommand：
%   \newcommand 遇到已存在的命令会直接报错
%       ! LaTeX Error: Command \DeckX already defined.
%   而 \renewcommand 又要求命令已存在 —— 单用哪一个都不行：
%   新组件用前者、覆盖旧定义用后者，可你事先不知道是哪种情况。
% \DeckDefine = \providecommand 后再 \renewcommand，两种情况都吃。
% 参数写法与 \newcommand 完全一致。
% ============================================================
\newcommand{\DeckDefine}[1]{%
  \providecommand{#1}{}%
  \renewcommand{#1}%
}

% ============================================================
% 二、页面骨架
% ============================================================

% --- 有标题页：标题由 main.tex 给出，正文来自页面文件 ---
\newcommand{\DeckSlide}[2]{%
  \begin{frame}[t]{#1}%
    \vspace*{\SlideBodyTopSkip}%
    \input{#2}%
  \end{frame}%
}

% --- 无标题页：封面与致谢页用 ---
% #1 内容文件，#2 垂直对齐（c 居中 / t 顶对齐）
% 做法是在这一页把 frametitle 模板清空（模板是局部的，不影响其他页），
% 内容用 \vfill 上下撑开，从而垂直居中。
\newcommand{\DeckBare}[2]{%
  \begin{frame}[#2]%
    \setbeamertemplate{frametitle}{}%
    \input{#1}%
  \end{frame}%
}

% ============================================================
% 三、最小原语（所有模板可以调用，它们自己不再依赖任何模板）
% ============================================================

% --- 正文行：全页只有这一种行 ---
% 姓名、单位、每一条要点，全部用同一个命令，因此字号、行距、
% 字重、颜色必然完全一致 —— 从代码层面就不可能出现不一致。
% 不要为"姓名"或"要点"另建命令：两个命令意味着两套样式。
% 悬挂缩进：折行时第二行与首行文字左对齐。
\newcommand{\DeckLine}[1]{%
  \par
  \hangindent=1.1em\hangafter=1
  \noindent #1\par
  \vspace*{\DeckRowSkip}%
}

% --- 分组留白 ---
% 唯一用来分组的手段。不用分隔线、不用小标题、不用字号差异。
\newcommand{\DeckGap}{%
  \par\vspace*{\DeckBlockSkip}%
}

% --- 超链接 ---
% beamer 已加载 hyperref，直接用 \href。
% 颜色用强调色，不加下划线边框（hyperref 的默认方框很难看）。
\newcommand{\DeckLink}[2]{%
  \href{#1}{\color{SlideAccent}#2}%
}

% --- 表格表头下的细线 ---
% 用 \noalign 单独发 \hrule，避免 \hline\noalign 给表格带来的异常深度
% （实测那一手会让 5 行表的深度变成 77pt）。
% T3 与 T4 都用它，所以归原语 —— 不能让 T4 依赖 T3。
\newcommand{\DeckCourseRule}{%
  \noalign{\vspace{0.09cm}\hrule height \DeckRuleHeight\vspace{0.13cm}}%
}
